\section{Related literature}\label{sec:literature}

The paper draws on three connected literatures: the run-theoretic literature on private near-money instruments, the microstructure literature on stablecoin arbitrage and secondary-market depth, and the payment-economics literature on cross-border corridors. The framing of the strategic stage builds on Morris-Shin global games \citep{morris_shin_2003_global_games} and on the bank-run formulation of \citet{goldstein_pauzner_2005_demand_deposits} that underpins \citet{diamond_dybvig_1983}.

\paragraph{Stablecoin runs as coordination games.} Several recent contributions have brought run-theoretic tools to stablecoins. \citet{ahmed_aldasoro_duley_2025} extend the Morris-Shin coordination problem along the information dimension by treating the volatility of reserve assets as unknown, which produces a second-generation global game with multiple locally-unique equilibria; their main object of study is the distinct effect of transparency and reserve quality. \citet{bertsch_2025_adoption_fragility} adopt a modified Goldstein-Pauzner formulation that combines an ex-ante adoption stage and an interim run game, and finds that most factors that raise the attractiveness of stablecoins also reduce their fragility, while wider adoption can be destabilising when the marginal holder is more flighty. \citet{gorton_klee_ross_ross_vardoulakis_2026} develop a leverage-money model in which speculative crypto demand for stablecoins generates the lending-rate premium that compensates run risk; they nest a Goldstein-Pauzner bank-run game inside a Gromb-Vayanos leveraged-trading block and apply the model to the May 2022 turmoil. Our model retains the global-game backbone of these papers but jointly carries $K$ heterogeneous holder classes and couples the strategic stage to a payment-rail block, neither of which is present in the cited works.

\paragraph{Two-layered market structure and arbitrage.} \citet{ma_zeng_zhang_2025} document that arbitrage on the largest stablecoins is highly concentrated, with very few authorised redeemers acting at the issuer per month, and show that more efficient arbitrage can raise run risk by reducing the secondary-market price impact that disciplines investor selling. \citet{lyons_viswanath_natraj_2023} characterise the role of arbitrageurs in restoring stablecoin pegs, with safer collateral associated with greater peg stability. \citet{watsky_etal_2024} document large divergences between primary and secondary markets for stablecoins under stress. \citet{du_sonawane_watsky_2025} provide a detailed reconstruction of the March 2023 USDC episode and treat it as a shadow bank run. Our dual-channel exit structure (primary redemption at the issuer plus secondary-market sale with finite depth) follows this strand; the difference is that the heterogeneity in our model resides among end users rather than between arbitrageurs and an undifferentiated investor mass, which lets us tie each class to a specific payment corridor.

\paragraph{Cross-border evidence and payment economics.} The paper's corridor mapping is informed by descriptive work on cross-border stablecoin flows. \citet{auer_lewrick_paulick_2024_defiying_gravity} document that stablecoin activity responds to international frictions more strongly than native cryptoassets do, with usage concentrated in corridors where traditional rails are costly. The corridor benchmarks themselves come from the World Bank Remittance Prices Worldwide quarterly panel for 2023-2024, which covers the US to Mexico and US to the Philippines corridors with sufficient depth to identify the per-unit incumbent cost. Fragmentation is a structural feature of the rails on which stablecoins circulate; \citet{shin_2026_tokenomics_fragmentation} and \citet{boissay_cornelli_doerr_frost_2022} argue that scalability constraints split liquidity across chains and venues, with implications for the law of one price and for coordination dynamics. Our payment-rail block treats per-class congestion as a parameter rather than as an endogenous object, which is enough to give the cost-benefit comparison empirical bite without taking a stand on the equilibrium structure of fees and validators.

\paragraph{Policy and regulation.} The cost-benefit exercise is sensitive to discrete policy events that re-anchor the equilibrium. The November 2025 package of resolutions issued by the Banco Central do Brasil frames virtual assets within the foreign-exchange market and the international-capital framework, and a follow-up resolution issued in 2026 confirms that stablecoins cannot be used in electronic foreign-exchange operations. The package is documented in the primary materials made public by the Banco Central do Brasil. \citet{arner_auer_frost_2020} review the broader regulatory landscape; the question of how regulation interacts with peg stability is also taken up in the run-theoretic work cited above.

\paragraph{Position of the paper.} The paper jointly carries (i) type-heterogeneous Morris-Shin thresholds, (ii) a payment-system congestion block coupled to the indifference condition, and (iii) a corridor-level certainty-equivalent cost comparison against named incumbent services. The first piece extends the run-theoretic literature into a setting where the cross-section of holder types matters for the equilibrium. The second piece embeds the strategic stage in a payment-economics object rather than a pure asset-pricing object. The third piece anchors the model to data that can be read by both finance and payment-policy audiences. Each piece on its own is incremental; the joint object is the contribution.
